% Unfinished round 237 draft. The cut replay was stopped after 300 of 623 statements.
% Update verification claims from a final receipt before using this passage.
\subsection{An exchange in the larger parity class}\label{sec:majority237}
A working hypothesis is that $a(9)=191$. A longer snake would have a
192-edge prefix. To look for such a prefix, let $B$ be either 96-vertex
parity class of a saved 191-edge path. Choose $d\in B$ and distinct
$a,b$ of the same parity outside $B$. Set
\[
 U=(B\setminus\{d\})\cup\{a,b\}.
\]
Can these 97 vertices occur at all even positions of a 192-edge snake?
The other 96 vertices, both endpoints and the ordering are free. This
question tests one consequence of the working hypothesis. The numerical
conjecture in Section~\ref{sec:hypothesis} makes a different prediction.

For any one-parity set $W$, form the simple graph $G(W)$ on $W$.
Join $v$ and $w$ if some opposite-parity vertex has exactly $v$ and $w$
as neighbours in $W$. Multiple such labels give one edge.
Section~\ref{sec:majority237} uses the following equivalence. A snake
with its even positions equal to $U$ orders $U$ along a Hamiltonian
path of $G(U)$, since every intervening odd vertex has exactly its two
path neighbours in $U$. Conversely, interleave a Hamiltonian order
with one label per edge. Distinct edges in the order require distinct
labels, and the two-neighbour condition excludes every chord.

\begin{lemma}\label{lem:partial237}
If $G(U)$ admits a Hamiltonian path, $W\subseteq U$, and $r=|U\setminus W|$,
then $2I(W)+L(W)\le 2r+2$, where $I(W)$ and $L(W)$ count the isolated
vertices and the vertices of degree one in $G(W)$.
\end{lemma}
\begin{proof}
Delete the $r$ missing positions from a Hamiltonian order. At most
$r+1$ nonempty runs remain. Every consecutive old pair within a run
retains its label when restricted to $W$. The runs thus form a spanning
union of paths in $G(W)$. In this union the sum of $2-\deg(v)$ is twice
the number of nonempty runs, including singleton runs. Additional edges
can only lower the sum of $\max(0,2-\deg(v))$, which in $G(W)$ is
$2I(W)+L(W)$.
\end{proof}

\begin{lemma}\label{lem:cut237}
Suppose $G(U)$ admits a Hamiltonian path. For arbitrary subsets
$S,T\subseteq U$, let $e(S,U\setminus T)$ count ordered adjacent pairs
$(u,v)$ with $u\in S$ and $v\in U\setminus T$. Then
\[
 2|S|\le 2|T|+2+e(S,U\setminus T).
\]
\end{lemma}
\begin{proof}
Sum the path degrees at the vertices in $S$. Only the two endpoints can
have degree one, so the sum is at least $2|S|-2$. Count these incidences
by their destinations. Each vertex of $T$ receives at most two; every
remaining incidence gives a different ordered edge counted by $e$.
This also covers overlapping $S$ and $T$.
\end{proof}

The 145 distinct saved classes give
$145\cdot96\binom{160}{2}=177\,062\,400$ descriptions before
deduplication. With $a<b$, removing $b$ from the proposed order requires
$2I+L\le4$ at $W=(B\setminus\{d\})\cup\{a\}$. This excludes
$170\,560\,947$ descriptions. At the final set, an isolated vertex or
more than two leaves excludes another $6\,500\,774$. The remaining
679 descriptions represent 623 distinct sets. A separate enumeration
applied the single-addition condition to both $a$ and $b$ and reconstructed
the surviving graphs directly. It returned exactly the same descriptions.

Every one of the 623 sets violates Lemma~\ref{lem:cut237} for explicit
$S$ and $T$. For example, take saved class 43, remove vertex 138 and add
8 and 158. With the integer cube labels of this paper, choose
\[
 S=\{56,88,158,188,280,440,446,472\},\qquad T=\{176,336\}.
\]
Exactly eight ordered edges leave $S$ for $U\setminus T$. The required
inequality becomes $16\le14$. Across all 623 sets, the excess of the
left side over the right side is 2, 4, 6 or 8, with respective
frequencies 161, 291, 167 and 4. Each count can be checked within the
cube: from a source $u$, enumerate the 36 unordered pairs of distinct
axes, giving targets $w$ at distance two. There are exactly two
intermediate cube vertices; either gives an edge if its selected
neighbours are exactly $u$ and $w$. Count the target once even if both
labels work. The finite enumeration was independently reproduced, and
all 623 displayed-by-data inequalities have individual kernel checks.
An independent necessary-degree search excludes the same 145 families
before imposing connectivity~\cite{research237}. The correspondence from
the complete enumeration to the saved bank is computational evidence.

\begin{corollary}
Under the finite covers above and those of the preceding three rounds,
a 192-edge snake shares at most 186 vertices with each saved 191-edge
snake or any cube-isometric copy.
\end{corollary}
\begin{proof}
The candidate has parity classes of sizes 97 and 96. The larger class
cannot retain 96 seed vertices by the fixed-class survey, or exactly 95
by the present exclusion. The smaller class retains at most 92 by the
earlier covers. Thus the intersection has size at most $94+92=186$.
The seed has 192 vertices, the candidate 193, giving at least six
deletions and seven additions. Applying an inverse cube isometry gives
the same conclusion for each copy.
\end{proof}

\subsection{Recovering a selection with two missing counts}\label{sec:erasure237}
\begin{proposition}
For any two distinct vertices of $\Q_9$, the neighbour counts at the
remaining 510 vertices uniquely determine a binary vertex selection.
\end{proposition}
\begin{proof}
Let $z$ be the difference of two selection indicators and let $b(v)$
be the sum of $z$ at the neighbours of $v$. Suppose $b$ vanishes off
$\{p,q\}$. Write $x=b(p)$ and $y=b(q)$. Apply the inverse from
Proposition~\ref{prop:inverse236} at the antipodes of $p$ and $q$.
Its distance-nine coefficient is 128. The two off-diagonal distances
are equal and below nine, so their common coefficient $c$ belongs to
$\{0,35,-10,8,-16\}$. For some $s,t\in\{-1,0,1\}$ this gives
\[
 128x+cy=315s,\qquad cx+128y=315t.
\]
If exactly one of $s,t$ is nonzero, the corresponding unknown has
absolute value $315\cdot128/(128^2-c^2)$, strictly between two and
three. This contradicts integrality. If both are nonzero, then $t=s$
or $t=-s$. Inverting the two equations gives
$|x|=315/(128+c)$ or $315/(128-c)$, respectively. The possible
denominators are $93,112,118,120,128,136,138,144,163$, none of which
divides 315. Hence $s=t=0$. Since $|c|<128$, we get $x=y=0$.
The complete profile now vanishes, and the inverse gives $z=0$.
\end{proof}
The inverse and this two-erasure result were checked on the full
nine-dimensional cube. The result concerns binary selections; the proof
uses the three possible values of their difference. Retained path
vertices alone generally leave many more than two counts unknown.

